\begin{table}[htbp]
\centering
\caption{Air Quality Threshold Exceedance}
\label{tab:thresholds}
\begin{tabular}{lrr}
\toprule
Threshold & N Readings & Percentage \\
\midrule
PM2.5 $>$ Good ($>$9 $\mu$g/m$^3$) & 98 & 100.0\% \\
PM2.5 $>$ Moderate ($>$35.4 $\mu$g/m$^3$) & 97 & 99.0\% \\
PM2.5 $>$ Unhealthy for Sensitive ($>$55.4 $\mu$g/m$^3$) & 97 & 99.0\% \\
PM2.5 $>$ Unhealthy ($>$125.4 $\mu$g/m$^3$) & 21 & 21.4\% \\
PM2.5 $>$ Very Unhealthy ($>$225.4 $\mu$g/m$^3$) & 6 & 6.1\% \\
\midrule
CO $>$ WHO 24-hour ($>$4000 $\mu$g/m$^3$) & 0 & 0.0\% \\
CO $>$ EPA 8-hour ($>$10305 $\mu$g/m$^3$) & 0 & 0.0\% \\
\bottomrule
\end{tabular}
\begin{tablenotes}
\small
\item Note: EPA 2024 lowered the annual PM2.5 standard from 12 to 9 $\mu$g/m$^3$.
\item CO readings well below WHO/EPA thresholds (max $\approx$1500 $\mu$g/m$^3$ $\approx$ 1.3 ppm).
\end{tablenotes}
\end{table}
